\documentclass[a4paper,12pt]{article}

\usepackage[utf8]{inputenc}
\usepackage[T1]{fontenc}
\usepackage{babel}
\usepackage{graphicx}
\usepackage{geometry}
\usepackage{array}
\usepackage{listings}
\usepackage{xcolor}
\usepackage{hyperref}
\usepackage{float}
\usepackage{tabularx}
\usepackage{booktabs}
\usepackage{textcomp}
\geometry{margin=2.5cm}

% Style pour le code
\definecolor{codebg}{rgb}{0.97,0.97,0.97}
\definecolor{codecomment}{rgb}{0.0,0.5,0.0}
\definecolor{codekw}{rgb}{0.0,0.0,0.7}
\definecolor{codestr}{rgb}{0.6,0.1,0.1}

\lstset{
    basicstyle=\ttfamily\footnotesize,
    breaklines=true,
    frame=single,
    backgroundcolor=\color{codebg},
    keywordstyle=\color{codekw}\bfseries,
    commentstyle=\color{codecomment},
    stringstyle=\color{codestr},
    numbers=left,
    numberstyle=\tiny\color{gray},
    numbersep=5pt,
    showstringspaces=false
}

\title{
    \vspace{-1cm}
    {\LARGE\bfseries Rapport de Projet}\\[0.2cm]
    {\large Mini Site E-commerce de Pizzeria}\\[0.1cm]
    {\Large\textbf{\color{red!70!black}AniPriZZA}}
}
\author{Judeken \& Peter \\ L2 Informatique}
\date{\today}

\begin{document}

\begin{titlepage}
\centering
\vspace*{2cm}
{\Huge\bfseries Rapport de Projet\par}
\vspace{0.5cm}
{\LARGE Mini Site E-commerce de Pizzeria\par}
\vspace{0.3cm}
{\Huge\color{red!70!black}\bfseries AniPriZZA\par}
\vspace{1cm}
\rule{\linewidth}{0.5mm}
\vspace{0.5cm}

\begin{tabular}{ll}
\textbf{Auteurs} & Judeken \& Peter canisious linton \\
\textbf{Niveau} & L2 Informatique \\
\textbf{Date} & 18/04/2026 \\
\textbf{Technologies} & PHP, MySQL, HTML/CSS/JavaScript \\
\end{tabular}

\vspace{0.5cm}
\rule{\linewidth}{0.5mm}

\vspace{1cm}
\begin{abstract}
Ce rapport présente la conception et la réalisation du projet \textbf{AniPriZZA}, un mini site e-commerce de pizzeria napolitaine développé en HTML/Css/Java/PHP/MySQL. Le système offre un catalogue de pizzas dynamique, un panier JavaScript, un système d'authentification sécurisé, un programme de fidélité, et une interface d'administration complète.
\end{abstract}
\end{titlepage}

\tableofcontents
\newpage

%--------------------------------------------------
\section{Introduction}

Dans ce projet, nous avons réalisé \textbf{AniPriZZA}, un mini site e-commerce de pizzeria napolitaine permettant de commander des pizzas en ligne.

L'objectif est de développer un site web dynamique complet en utilisant :
\begin{itemize}
    \item \textbf{PHP} --- logique serveur (sessions, traitement des formulaires, CRUD complet)
    \item \textbf{MySQL} --- persistance des données (pizzas, utilisateurs, commandes)
    \item \textbf{HTML / CSS} --- structure et design de l'interface (dark theme premium)
    \item \textbf{JavaScript} --- interactions client (panier dynamique, filtres, animations, LocalStorage)
\end{itemize}

Le projet intègre une architecture en couches avec séparation des responsabilités entre la présentation (HTML/CSS/JS), la logique métier (PHP) et les données (MySQL), et distingue trois niveaux d'accès : visiteur, client authentifié et administrateur.

%--------------------------------------------------
\section{Présentation du site}

\textbf{AniPriZZA} est une pizzeria napolitaine artisanale fondée en 2022. Le site propose une expérience e-commerce complète avec les pages suivantes :

\begin{itemize}
    \item \textbf{index.php} --- Page d'accueil : hero section (four à bois 450°C, cuisson 90s), pizzas phares statiques, caractéristiques de la maison
    \item \textbf{produits.php} --- Catalogue filtrable par catégorie (Classique, Spicy, Veggie, Premium), alimenté depuis la base de données
    \item \textbf{panier.php} --- Gestion dynamique du panier (JS + LocalStorage), validation de commande, système de points fidélité
    \item \textbf{login.php} --- Inscription, connexion, gestion du profil (téléphone, adresse), suppression de compte
    \item \textbf{contact.php} --- Formulaire de contact avec informations de la pizzeria
    \item \textbf{admin\_pizzas.php} --- Back-office : gestion complète du menu (CRUD + toggle disponibilité)
    \item \textbf{admin\_users.php} --- Back-office : gestion des comptes utilisateurs
\end{itemize}

\subsection{Acteurs du système}

Le système distingue trois types d'acteurs :

\begin{center}
\renewcommand{\arraystretch}{1.4}
\begin{tabular}{|>{\bfseries}p{2.5cm}|p{10cm}|}
\hline
Acteur & Description et droits \\
\hline
Visiteur & Consulter l'accueil, parcourir le catalogue, envoyer un message de contact. Pas d'authentification requise. \\
\hline
Client & Hérite des droits visiteur. Peut s'inscrire, se connecter, gérer son panier, passer des commandes, utiliser ses points fidélité, et gérer son profil. \\
\hline
Admin & Accès complet au back-office : CRUD sur les pizzas, toggle de disponibilité, gestion des comptes clients (ajout, suppression). \\
\hline
\end{tabular}
\end{center}

%--------------------------------------------------
\section{Fonctionnalités}

\subsection{Fonctionnalités côté client}

\begin{itemize}
    \item Affichage du catalogue de pizzas depuis la base de données (\texttt{SELECT} avec filtre \texttt{disponible = 1})
    \item Filtrage dynamique du catalogue par catégorie (JavaScript côté client, sans rechargement)
    \item Gestion du panier : ajout, modification de quantité, suppression (LocalStorage + JavaScript)
    \item Inscription avec 50 points de bienvenue offerts automatiquement
    \item Connexion sécurisée avec hash bcrypt (\texttt{password\_hash()} / \texttt{password\_verify()})
    \item Validation de commande avec calcul du total et des points gagnés (+2 pts/pizza)
    \item Utilisation des points de fidélité (100 pts = remise de 12,90\texteuro{} sur la commande)
    \item Mise à jour du profil utilisateur (téléphone, adresse de livraison)
    \item Suppression définitive du compte avec nettoyage de session
\end{itemize}

\subsection{Fonctionnalités côté administrateur}

\begin{itemize}
    \item Ajout d'une pizza (\texttt{INSERT} : nom, ingrédients, prix, image URL, catégorie)
    \item Modification d'une pizza (\texttt{UPDATE} via formulaire inline dans le tableau)
    \item Suppression d'une pizza (\texttt{DELETE} avec dialogue de confirmation)
    \item Activation/désactivation d'une pizza (toggle \texttt{disponible} 0/1) sans suppression
    \item Création manuelle d'un compte utilisateur (client ou administrateur)
    \item Suppression d'un compte utilisateur (avec protection anti-auto-suppression)
    \item Visualisation de tous les inscrits avec leurs points fidélité et leur rôle
\end{itemize}

%--------------------------------------------------
\section{Conception}

\subsection{Diagramme de cas d'utilisation}

Le diagramme ci-dessous représente les interactions entre les trois acteurs du système et les fonctionnalités disponibles.

\begin{figure}[H]
    \centering
    \includegraphics[width=\textwidth]{usecase.png}
    \caption{Diagramme de cas d'utilisation --- AniPriZZA}
    \label{fig:usecase}
\end{figure}

%--------------------------------------------------
\subsection{Diagramme de classes}

Le modèle de données du projet repose sur trois tables principales : \textbf{pizzas}, \textbf{users} et \textbf{messages}.

\begin{figure}[H]
    \centering
    \includegraphics[width=\textwidth]{classe.png}
    \caption{Diagramme de classes --- AniPriZZA}
    \label{fig:class}
\end{figure}

%--------------------------------------------------
\subsection{Diagramme de séquence --- Inscription d'un client}

Le scénario d'inscription illustre la validation des données et la création du compte.

\begin{figure}[H]
    \centering
    \includegraphics[width=\textwidth]{sequence_inscription.png}
    \caption{Diagramme de séquence --- Inscription d'un client}
    \label{fig:seq_inscription}
\end{figure}

\subsection{Diagramme de séquence --- Passage d'une commande}

Le scénario de commande illustre la logique de fidélité et la mise à jour des points.

\begin{figure}[H]
    \centering
    \includegraphics[width=\textwidth]{sequence_commande.png}
    \caption{Diagramme de séquence --- Passage d'une commande}
    \label{fig:seq_commande}
\end{figure}

%--------------------------------------------------
\section{Base de données}

Nous avons créé une base de données appelée \textbf{aniprizza} contenant trois tables principales.

\subsection{Table pizzas}

\begin{center}
\renewcommand{\arraystretch}{1.4}
\begin{tabular}{|l|l|l|}
\hline
\textbf{Champ} & \textbf{Type} & \textbf{Description} \\
\hline
id & INT AUTO\_INCREMENT PK & Identifiant unique \\
nom & VARCHAR(100) NOT NULL & Nom de la pizza \\
ingredients & TEXT NOT NULL & Liste des ingrédients \\
prix & DECIMAL(5,2) NOT NULL & Prix en euros \\
image\_url & TEXT & URL de l'image (Unsplash) \\
categorie & VARCHAR(50) DEFAULT 'classique' & Catégorie (classique/spicy/veggie/premium) \\
disponible & TINYINT(1) DEFAULT 1 & Disponibilité (0=masqué, 1=affiché) \\
\hline
\end{tabular}
\end{center}

\subsection{Table users}

\begin{center}
\renewcommand{\arraystretch}{1.4}
\begin{tabular}{|l|l|l|}
\hline
\textbf{Champ} & \textbf{Type} & \textbf{Description} \\
\hline
id & INT AUTO\_INCREMENT PK & Identifiant unique \\
nom & VARCHAR(100) NOT NULL & Nom ou pseudo \\
email & VARCHAR(150) NOT NULL UNIQUE & Adresse e-mail (clé unique) \\
password & VARCHAR(255) NOT NULL & Mot de passe hashé (bcrypt) \\
role & ENUM('user','admin') & Rôle (client ou administrateur) \\
points\_fidelite & INT DEFAULT 0 & Points cumulés (50 à l'inscription) \\
telephone & VARCHAR(20) & Téléphone (optionnel) \\
adresse & TEXT & Adresse de livraison (optionnel) \\
created\_at & TIMESTAMP DEFAULT NOW() & Date d'inscription automatique \\
\hline
\end{tabular}
\end{center}

\subsection{Table messages}

\begin{center}
\renewcommand{\arraystretch}{1.4}
\begin{tabular}{|l|l|l|}
\hline
\textbf{Champ} & \textbf{Type} & \textbf{Description} \\
\hline
id & INT AUTO\_INCREMENT PK & Identifiant unique \\
nom & VARCHAR(100) NOT NULL & Nom de l'expéditeur \\
email & VARCHAR(150) NOT NULL & Email de l'expéditeur \\
sujet & VARCHAR(255) NOT NULL & Sujet du message \\
message & TEXT NOT NULL & Contenu du message \\
created\_at & TIMESTAMP DEFAULT NOW() & Date d'envoi automatique \\
\hline
\end{tabular}
\end{center}

\subsection{Script SQL de création}

\begin{lstlisting}[language=SQL, caption=Script de création de la base de données AniPriZZA]
CREATE DATABASE IF NOT EXISTS aniprizza CHARACTER SET utf8mb4;
USE aniprizza;

CREATE TABLE pizzas (
    id          INT AUTO_INCREMENT PRIMARY KEY,
    nom         VARCHAR(100)         NOT NULL,
    ingredients TEXT                 NOT NULL,
    prix        DECIMAL(5,2)         NOT NULL,
    image_url   TEXT,
    categorie   VARCHAR(50)          DEFAULT 'classique',
    disponible  TINYINT(1)           DEFAULT 1
);

CREATE TABLE users (
    id               INT AUTO_INCREMENT PRIMARY KEY,
    nom              VARCHAR(100)     NOT NULL,
    email            VARCHAR(150)     NOT NULL UNIQUE,
    password         VARCHAR(255)     NOT NULL,
    role             ENUM('user','admin') DEFAULT 'user',
    points_fidelite  INT              DEFAULT 0,
    telephone        VARCHAR(20),
    adresse          TEXT,
    created_at       TIMESTAMP        DEFAULT CURRENT_TIMESTAMP
);

CREATE TABLE messages (
    id         INT AUTO_INCREMENT PRIMARY KEY,
    nom        VARCHAR(100)     NOT NULL,
    email      VARCHAR(150)     NOT NULL,
    sujet      VARCHAR(255)     NOT NULL,
    message    TEXT             NOT NULL,
    created_at TIMESTAMP        DEFAULT CURRENT_TIMESTAMP
);

\end{lstlisting}

%--------------------------------------------------
\section{Réalisation}

\subsection{Structure du projet}

\begin{verbatim}
/AniPriZZA/
|-- config.php           → Connexion BDD (MySQLi) + démarrage session
|-- index.php            → Page d'accueil (statique + dynamique)
|-- css/
|   \-- style.css        → Design dark premium (variables CSS, responsive)
|-- js/
|   \-- script.js        → Panier (LocalStorage), filtres, animations IntersectionObserver
\-- pages/
    |-- produits.php     → Catalogue dynamique (SELECT BDD, filtres JS)
    |-- panier.php       → Validation commande + logique fidélité (PHP)
    |-- login.php        → Inscription / Connexion / Profil / Suppression compte
    |-- contact.php      → Formulaire de contact (INSERT messages)
    |-- admin_pizzas.php → Back-office : CRUD pizzas + toggle disponibilité
    \-- admin_users.php  → Back-office : gestion utilisateurs (protégé admin)
\end{verbatim}

\subsection{Connexion à la base de données}

\begin{lstlisting}[language=PHP, caption=config.php — Connexion centralisée]
<?php
$host     = "localhost";
$user     = "root";
$password = "";          // XAMPP : vide par defaut
$dbname   = "aniprizza";

$conn = new mysqli($host, $user, $password, $dbname);
$conn->set_charset("utf8mb4");  // Support des emojis et accents

if ($conn->connect_error) {
    die(json_encode([
        "error"   => true,
        "message" => "Connexion echouee : " . $conn->connect_error
    ]));
}

// Demarrage session centralisee (evite les appels multiples)
if (session_status() === PHP_SESSION_NONE) {
    session_start();
}
?>
\end{lstlisting}

\subsection{Exemple : Ajout d'une pizza (admin)}

\begin{lstlisting}[language=PHP, caption=admin\_pizzas.php — Ajout d'une pizza (extrait)]
if ($_SERVER["REQUEST_METHOD"] === "POST" && isset($_POST['add'])) {
    $nom  = $conn->real_escape_string(trim($_POST['nom']));
    $ing  = $conn->real_escape_string(trim($_POST['ing']));
    $prix = (float)$_POST['prix'];
    $img  = $conn->real_escape_string(trim($_POST['img']));
    $cat  = $conn->real_escape_string(trim($_POST['categorie']));

    if ($nom && $ing && $prix > 0) {
        $conn->query("INSERT INTO pizzas 
                     (nom, ingredients, prix, image_url, categorie)
                     VALUES ('$nom','$ing',$prix,'$img','$cat')");
        $flash = "\"$nom\" ajoutee au menu !";
    }
}
\end{lstlisting}

\subsection{Système de fidélité}

Le programme de fidélité est géré entièrement côté serveur dans \texttt{panier.php} :

\begin{lstlisting}[language=PHP, caption=panier.php — Logique de fidélité (extrait)]
$uid = (int)$_SESSION['user_id'];
$total_qty   = (int)$_POST['total_qty'];
$total_price = (float)$_POST['total_price'];

// Recuperation des points actuels
$user_query = $conn->query(
    "SELECT points_fidelite FROM users WHERE id=$uid"
);
$points  = (int)$user_query->fetch_assoc()['points_fidelite'];
$discount = 0;

// Utilisation des points (100 pts = -12.90 euro)
if (isset($_POST['use_points']) && $points >= 100) {
    $discount = 12.90;
    $conn->query(
        "UPDATE users SET points_fidelite = points_fidelite - 100 
         WHERE id=$uid"
    );
}

// Calcul du total final
$final_price = max(0, $total_price - $discount);

// Credit des nouveaux points (+2 par pizza)
$points_earned = $total_qty * 2;
$conn->query(
    "UPDATE users SET points_fidelite = points_fidelite + $points_earned 
     WHERE id=$uid"
);
\end{lstlisting}

\subsection{Sécurité}

\begin{itemize}
    \item \textbf{Hashage des mots de passe} : algorithme bcrypt via \texttt{password\_hash(PASSWORD\_BCRYPT)} --- résistant aux attaques par table arc-en-ciel
    \item \textbf{Échappement des entrées} : utilisation systématique de \texttt{real\_escape\_string()} pour toutes les variables insérées en SQL
    \item \textbf{Protection des routes admin} : vérification de \texttt{\$\_SESSION['user\_role'] === 'admin'} avant tout accès au back-office, avec redirection vers login.php
    \item \textbf{Protection anti-auto-suppression} : vérification que l'ID à supprimer est différent de l'ID de la session admin courante
    \item \textbf{Cast des identifiants} : tous les IDs provenant de l'URL ou du formulaire sont castés en \texttt{int} pour prévenir les injections
\end{itemize}

%--------------------------------------------------
\section{Interface du site}

\subsection{Design}

Le site adopte un design \textit{dark premium} inspiré des grandes pizzerias italiennes :

\begin{itemize}
    \item Fond sombre (\texttt{\#0f0f0f}) avec accents rouge-orangé (\texttt{\#e5641e}) et or (\texttt{\#c9a96e})
    \item Typographie mixte : \textit{Cormorant Garamond} (display) + \textit{Inter} (corps de texte)
    \item Variables CSS pour la cohérence des couleurs sur toutes les pages
    \item Animations d'apparition (\texttt{slide-up}) déclenchées par \texttt{IntersectionObserver}
    \item Design entièrement responsive avec grille CSS adaptative et hamburger menu mobile
\end{itemize}

\subsection{Pages principales}

\begin{center}
\renewcommand{\arraystretch}{1.4}
\begin{tabular}{|l|p{8.5cm}|}
\hline
\textbf{Page} & \textbf{Description} \\
\hline
Accueil & Hero section, stats (450°C, 90s, 100\% artisanal), 3 pizzas phares, section features \\
La Carte & Grille de pizzas filtrables, boutons «Ajouter» au panier, badge catégorie coloré \\
Panier & Récapitulatif dynamique, validation, case «utiliser mes points», message de confirmation \\
Mon Compte & Formulaires inscription/connexion, tableau de bord client (points, coordonnées) \\
Admin Pizzas & Tableau inventaire complet, formulaire d'ajout, boutons modification/suppression/toggle \\
Admin Users & Liste des inscrits avec rôle, points, date; formulaire de création manuelle \\
\hline
\end{tabular}
\end{center}

%--------------------------------------------------
\section{Difficultés rencontrées}

Au cours du développement, nous avons rencontré plusieurs difficultés techniques et fonctionnelles :

\begin{itemize}
    \item \textbf{Gestion des sessions PHP} : assurer la cohérence de l'état de connexion entre toutes les pages, notamment pour l'affichage conditionnel des menus (espace client vs visiteur) et la protection des routes admin par redirection
    
    \item \textbf{Architecture du panier} : le choix entre un panier en LocalStorage (côté client, rapide et sans rechargement) et un panier en base de données (persistant, lié au compte) a été tranché en faveur du LocalStorage, avec validation des données côté serveur uniquement au moment du checkout
    
    \item \textbf{Sécurité des formulaires} : prévention des injections SQL via l'échappement systématique des entrées, la validation des types (cast explicite des IDs en entier), et la vérification des conditions d'accès
    
    \item \textbf{Toggle de disponibilité} : implémentation du switch on/off en CSS pur avec redirection PHP après mise à jour, en évitant l'usage d'un framework JavaScript lourd
    
    
    \item \textbf{Calcul du système de fidélité} : gestion des cas limites (total inférieur à la réduction, protection contre les points négatifs, scénario où l'utilisateur n'est pas connecté au moment du checkout)
\end{itemize}

%--------------------------------------------------
\section{Conclusion}

Ce projet nous a permis de comprendre et d'appliquer :

\begin{itemize}
    \item La conception complète d'une application web (diagrammes UML : cas d'utilisation, classes, séquence)
    \item Le développement full-stack d'un site e-commerce dynamique avec PHP et MySQL
    \item Les opérations CRUD complètes sur plusieurs entités (pizzas et utilisateurs)
    \item La gestion sécurisée des sessions et de l'authentification (bcrypt, vérification de rôle)
    \item L'implémentation d'une logique métier complexe (système de points fidélité)
    \item La séparation des responsabilités et la protection des ressources sensibles
\end{itemize}

\subsection{Améliorations possibles}

\begin{itemize}
    \item Historique des commandes persistant en base de données avec table \texttt{commandes} et \texttt{commande\_items}
    \item Module d'upload d'images pour les pizzas (au lieu des URL externes Unsplash)
    \item Suivi de commande en temps réel avec statuts (en préparation, en livraison, livrée)
    \item Intégration d'un système de paiement sécurisé (Stripe, PayPal)
    \item Système d'avis et de notation des pizzas (table \texttt{avis}, calcul moyenne)
    \item Développement d'une API REST pour une application mobile (React Native, Flutter)

%--------------------------------------------------
\section{Bibliographie}

\begin{itemize}
    \item Les différents chapitres fournis en cours 
    \item Documentation MySQL : \url{https://dev.mysql.com/doc/}
    \item MDN Web Docs (HTML/CSS/JavaScript) : \url{https://developer.mozilla.org/fr/}
    \item Cours de Conception d'Applications Web
    \item Exemple de Conception Application Web --- Document fourni en cours
    \item UML 2.5 --- Object Management Group : \url{https://www.omg.org/spec/UML/}
\end{itemize}

\end{document}